\subsection{Motivation}
\label{sec:intro:motivation}

Large language models (LLMs) are increasingly deployed not as isolated
predictors but as interacting agents that reason, argue, and decide in groups.
A prominent instance of this shift is \emph{multi-agent debate} (MAD), in which
several LLM instances independently propose answers, exchange their reasoning
over multiple rounds, and revise their positions before the group commits to a
final decision \citep{du2024improving}. Across factual and mathematical
reasoning tasks, this simple protocol has been reported to improve accuracy and
reduce hallucination relative to a single model answering alone
\citep{du2024improving, liang2024encouraging, chen2024reconcile}, and it now
underpins a growing family of collaborative and evaluative systems
\citep{chan2024chateval}. The appeal is intuitive: independent agents make
partially independent errors, and deliberation offers a mechanism for the group
to identify and correct them.

This optimistic reading, however, rests almost entirely on \emph{outcome}
measurements. The literature evaluates multi-agent systems predominantly by
final task performance, treating the deliberation itself as an opaque process
that either does or does not yield a better answer. Recent critical
reassessments show that this outcome-centric view is fragile: debate frequently
fails to outperform single-model baselines, collapses under controlled comparison
with simpler alternatives, and exhibits failure rates between 41\% and 86.7\%
across seven frameworks \citep{liu2025stopovervaluing, zhang2025debateorvote,
bertalanic2026costofconsensus, cemri2025mast}. More
troublingly, the failures that do occur are not random noise but exhibit
structure. When two LLMs debate, both tend to grow \emph{more} confident that
they have won as the exchange proceeds, with confidence escalating away from a
rational baseline even when no new supporting evidence is introduced
\citep{prasad2025certain}. Under distributed information, groups of LLM agents
converge prematurely on the evidence they already share and fail to surface the
private knowledge needed for the correct answer, so that collective accuracy
collapses far below what a single agent achieves with complete information
\citep{li2026hiddenbench}. Purely cooperative framings can further collapse the
debate into a mutually agreeing echo chamber, a failure mode termed
Degeneration-of-Thought \citep{liang2024encouraging}.

What these pathologies have in common is that they are \emph{dynamical}. Rising
confidence, premature lock-in, and conformity cascades are properties of how
positions evolve and influence one another across rounds, not of any single
agent's competence at a single point in time. An outcome metric, by
construction, cannot see them: two systems can reach the same final accuracy
through entirely different trajectories, one through genuine error correction
and the other through a confident slide into a shared mistake. Understanding
\emph{when} and \emph{why} multi-agent deliberation succeeds therefore requires
a vocabulary for the interaction dynamics themselves. \citet{miehling2025agentic}
make precisely this argument at a conceptual level, contending that agentic AI
must be studied through a systems-theoretic lens because emergent group-level
behaviour cannot be predicted from individual agent properties alone. Their
position, however, remains programmatic and offers no empirical instrument for
detecting such behaviour in practice.

Complex-systems theory supplies exactly such an instrument. Systems of many
interacting units that produce ordered collective behaviour without central
control, the defining feature of self-organization \citep{ashby1947principles},
are known to fall into a small number of recurring dynamical regimes. These
regimes have precise names and precise signatures: \emph{self-reinforcement},
the positive feedback by which a state amplifies its own persistence
\citep{arthur1989competing}; \emph{limit cycles}, self-sustained periodic
oscillations that never settle \citep{strogatz2015nonlinear};
\emph{bistability}, the coexistence of multiple stable outcomes selected by
initial conditions \citep{feudel2008complex}; and \emph{dominance}, the
concentration of influence onto a single node \citep{degroot1974consensus,
jia2015opinion}. Such patterns are studied as \emph{motifs}, recurring structural or
dynamical building blocks whose presence in a system is both detectable and
diagnostic of its behaviour \citep{milo2002network, alon2007network,
driscoll2024dynamical}. The pathologies documented in the MAD literature read
naturally as instances of this vocabulary: escalating confidence is a
self-reinforcement signature, premature consensus is a fixed-point attractor
reached through basin selection, and herding to a confident leader is a
dominance motif. The connection has begun to be made.
\citet{ashery2025conventions} demonstrate the spontaneous emergence of
universally adopted social conventions in decentralised LLM populations, show
that strong collective biases arise even when individual agents exhibit none,
and find that committed adversarial minorities can impose alternative
conventions on the majority. \citet{okawa2026biasedconsensus} derive a formal model predicting that once
conformity pressure among agents exceeds a critical threshold, the group is
pulled toward a shared bias regardless of what any individual agent believes.
What these treatments share is a substrate in which no option is better than any
other, so they can show that a group locks in but not whether it locks in on the
right answer; \citet{mehdizadeh2026topologymemory} name the epistemic case, where
some claims are better supported than others, as a critical open question.

This thesis addresses that gap. We construct a controlled cooperative
multi-agent debate system in which homogeneous LLM agents deliberate over
multiple rounds on tasks with verifiable ground truth, and we develop a suite of
detectors, each grounded in the complex-systems and opinion-dynamics literature,
that identify self-reinforcement, limit cycles, bistability, and dominance in
the resulting interaction logs. We run this system across a controlled grid of
communication topologies and memory configurations on two complementary
benchmarks: GPQA-Diamond \citep{rein2024gpqa}, a graduate-level science
question-answering benchmark where agents typically start with reasonable priors,
and HiddenBench \citep{li2026hiddenbench}, a hidden-profile benchmark
constructed so that agents start systematically wrong and success requires
genuine information pooling. The contrast between the two isolates whether
observed dynamics depend on the presence of a reliable individual signal or
persist even when the group must correct a shared bias. Throughout, we treat
motif detection not as a competitor to accuracy but as an explanatory layer
beneath it: a way to characterise the trajectory that produces an outcome, and
thereby to understand which design choices shape the deliberative dynamics of
LLM collectives.

This explanatory layer has a practical purpose. As multi-agent systems are
deployed in consequential settings, an operator needs to know not only whether a
configuration tends to produce correct answers but whether a given deliberation
is proceeding soundly or pathologically, whether the group is genuinely
correcting an error or confidently converging on a shared mistake, and whether it
will terminate at all. Outcome metrics cannot answer these questions.
Instruments are beginning to appear from other directions: online monitors
based on conditional transfer entropy over execution traces
\citep{caspian2026} and learned classifiers that predict from interpretable
features whether a query warrants debate at all \citep{fan2026imad}. What does
not yet exist is an instrument grounded in the dynamical vocabulary itself, one
that names which regime a deliberation is in rather than only flagging that
something is wrong. A
central thread of this thesis is therefore that each motif detector, being
computed from ordinary interaction logs, doubles as a candidate \emph{diagnostic}
of deliberative health. We do not deliver a validated diagnostic tool; we deliver
the empirical groundwork for one, a set of measurable signatures, characterised
across a controlled grid, together with tiered, evidence-graded design principles
that state honestly how far each is supported by our data.

In one sentence, the objective of this thesis is to understand \emph{how} a
cooperative LLM debate arrives at its answer, by measuring the self-organization
patterns in its trajectory, and to turn that understanding into evidence-graded
guidance for building such systems.


\subsection{Research Questions}
\label{sec:intro:rqs}

The overarching goal of this thesis is to establish whether the dynamical
vocabulary of complex systems can be applied empirically to LLM-based
multi-agent debate, and to characterise the conditions under which
self-organization patterns emerge. We decompose this goal into three research
questions, ordered by the class of evidence each requires.

\begin{itemize}
  \item[\textbf{RQ1}] \textbf{Measurement.} Can self-organization motifs,
        specifically self-reinforcement, limit cycles, bistability, and
        dominance, be detected in cooperative LLM-based multi-agent debate using
        detectors grounded in the complex-systems and opinion-dynamics literature
        and each tested against a principled null model? And which of these motifs
        actually occur, at what prevalence, and with what structure?

  \item[\textbf{RQ2}] \textbf{Characterisation.} What governs the presence and
        strength of each motif, and what do they predict about system outcomes?
        We examine the two canonical structural axes, communication topology and
        memory window, together with the information regime of the task and the
        composition of the opening votes, and ask what each motif implies for
        accuracy, convergence speed, and worst-case runtime.

  \item[\textbf{RQ3}] \textbf{Causation.} Is the composition of the opening votes
        causally responsible for the outcome? Can the system generate a correct
        answer that no agent initially held, or does deliberation function
        primarily as an amplifier of the initial vote distribution?
\end{itemize}

One commitment underpins these questions. Every detector is tested against a null model that destroys the pattern we look for and leaves everything else intact. This matters because our state space is small: with four agents and three or four options, votes repeat and influence concentrates by chance alone. The null tells us what chance looks like, so a near-zero prevalence becomes a measurement rather than a failure. We report the motifs we do not find alongside those we do.


\subsection{Contributions \& Structure}
\label{sec:intro:contributions}

This thesis makes three contributions, one for each research question.

\begin{itemize}
  \item \textbf{An instrument for measuring debate dynamics.} We define a
        two-phase synchronous debate protocol that separates public messages
        from private beliefs and confidences, so that the interaction log
        exposes exactly the quantities motif detection requires
        (\autoref{sec:exp}). On top of it we specify \suitename{}
        (Self-Organization Motif Analysis): detectors for self-reinforcement,
        bistability, dominance, and limit cycles, each grounded in the
        complex-systems or opinion-dynamics literature and each tested against
        its own null model (\autoref{sec:method}). Every detector states what it
        can and cannot claim, distinguishing a measured signature from a
        certified attractor. Running $R = 50$ repetitions of every task--config
        cell is what makes across-repetition attractor statistics measurable at
        all.

  \item \textbf{An empirical characterisation at scale.} We apply \suitename{}
        to 30{,}000 debates across twelve configurations and two deliberately
        contrasting benchmarks: a refinement regime in which agents mostly start
        correct, and a correction regime in which they start systematically
        wrong (\autoref{sec:results}). Self-reinforcement is universal and drives
        convergence, bistability is common and is a stable property of the task
        rather than of the configuration, dominance emerges behaviourally even
        under a symmetric topology, and genuine limit cycles are essentially
        absent outside the star, where they arise through a hub-driven
        mechanism. Reading the emergent influence graph within fully-connected
        runs reveals a lever the structural topology parameter cannot reach: a
        rare effective star, formed preferentially from diverse opening
        positions, is markedly more accurate than the fully-connected mean.
        Throughout, the structural axes we varied move the dynamics but not the
        accuracy.

    \item \textbf{A causal test of the initial-condition effect.} A
        seeded-initialisation experiment manipulates only the composition of the
        opening votes, leaving the debate itself unchanged. Inserting a single
        correct agent into an otherwise-incorrect opening raises accuracy from
        $1.8\,\%$ to $23.7\,\%$, and does so in all 34 tasks tested. In the
        unseeded arm, where no agent starts correct, the system arrives at the
        correct answer in under $2\,\%$ of repetitions, and its ability to do so
        declines as tasks get harder. A devil's advocate probe confirms the
        complementary vulnerability under a controlled composition: a single
        agent that never updates its position outweighs two agents that start
        correct, collapsing accuracy from $63.7\,\%$ to $10.0\,\%$. Together
        these upgrade the central claim from association to intervention: within
        the setting we tested, deliberation functions predominantly as a
        selective amplifier of its opening vote distribution
        (\autoref{sec:results}).
\end{itemize}

\noindent From these results we derive design guidance graded by how strongly our
data support each claim, and show that every \suitename{} detector, computed from
ordinary interaction logs, is a candidate observable for monitoring deliberative
health (\autoref{sec:discussion}). We deliver the measurements and the mapping
from principle to observable, not a packaged tool.

\paragraph{Structure of this thesis.}
The thesis follows the arc of the three research questions. \emph{Background}
(\autoref{sec:infsub}) introduces the vocabulary of self-organization and
dynamical motifs and fixes what we mean by an LLM agent and a debate;
\emph{Related Work} (\autoref{sec:related}) surveys how multi-agent systems are
currently analysed and locates the gap this thesis addresses. The instrument is
then built in two steps: \emph{System and Experimental Setup}
(\autoref{sec:exp}) specifies the debate protocol, the two contrasting
benchmarks, and the experimental grid, and \emph{Detection Methodology}
(\autoref{sec:method}) defines the four \suitename{} detectors together with the
null model each is tested against. \emph{Results} (\autoref{sec:results}) applies
the instrument, reporting accuracy and efficiency, the four motifs in turn, the
emergent influence topology, and the seeded-initialisation experiment.
\emph{Discussion} (\autoref{sec:discussion}) reads the motifs together as one
amplifier mechanism and distils the findings into evidence-graded design
principles. \emph{Limitations and Future Work} (\autoref{sec:limitations}) states
the scope of the claims and the agenda they imply, and
\emph{Conclusion} (\autoref{sec:conclusion}) returns to the three research
questions.